%% Finitary normal modal logics and descriptive general-frame completeness.

\subsection{Normal logics for arbitrary finitary operator signatures}
\label{sec:p10-polyadic-modal-algebras}

The algebraic and local completeness arguments extend to the full
signature of Definition~\ref{def:p10-bao}.  The number of operator
symbols need not be finite; each individual operator has finite arity.
The following language and normality conventions follow
\cite[preprint \S2, pp.~4, 6]{Venema2005Preprint}.

\begin{hbdefinition}{def:p10-polyadic-normal-logic}{Finitary normal modal language and logic}
Fix a set \(\Lambda\) with an arity \(n_f\in\mathbb N\) for each
\(f\in\Lambda\), and a countably infinite set \(X\) of proposition
variables disjoint from the connective symbols.
The set \(\operatorname{Fm}_\Lambda(X)\) consists of the finite formulas
\[
 \varphi ::= p\ (p\in X)\mid\bot\mid\neg\varphi
       \mid(\varphi\wedge\psi)\mid f(\varphi_1,\ldots,\varphi_{n_f}).
\]
The other Boolean connectives are abbreviations.
For \(n_f=0\), the expression \(f()\) is a formula constant.
Substitution replaces variables simultaneously by formulas and leaves
every connective, including a nullary one, fixed.

For \(n_f>0\), fix a coordinate \(1\leq i\leq n_f\) and formulas in the other
coordinates, and write \(f_i(\chi)\) for the resulting formula with
\(\chi\) in coordinate \(i\).
A \emph{normal \(\Lambda\)-logic} is a set
\(L\subseteq\operatorname{Fm}_\Lambda(X)\) containing all classical
propositional tautology instances and, for every such context, the axioms
\[
 \neg f_i(\bot),\qquad
 f_i(\varphi\vee\psi)\to
                 \bigl(f_i(\varphi)\vee f_i(\psi)\bigr).
\]
It is closed under modus ponens and uniform substitution.
It also has \emph{theorem-level coordinate monotonicity}: whenever
\(\varphi\to\psi\) belongs to \(L\), so does
\(f_i(\varphi)\to f_i(\psi)\), for every choice of the other arguments.
Write \(\vdash_L\varphi\) for \(\varphi\in L\), and call \(L\)
consistent when \(\bot\notin L\).

The reverse additivity implication follows from monotonicity applied
to \(\varphi\to\varphi\vee\psi\) and
\(\psi\to\varphi\vee\psi\).
The source states monotonicity for simultaneous changes of all arguments;
the coordinate form is equivalent by finite iteration, using reflexive
implications for unchanged arguments.  Its normality axiom puts
\(\top\) in the other coordinates.  Monotonicity in those coordinates
gives the arbitrary-context version displayed here, and the converse is
an instance of that version.
No coordinate rule or normality axiom is imposed when \(n_f=0\).
In particular, normality does not assert \(f()\leftrightarrow\bot\).

A derivation from a premise set \(\Gamma\) is a finite sequence using
members of \(\Gamma\), members of \(L\), and modus ponens.
Write \(\Gamma\vdash_L\varphi\) when such a sequence ends in
\(\varphi\).  Substitution and coordinate monotonicity apply to
theorems of \(L\), not freely to premise-dependent lines.
For one unary symbol \(\Diamond\), putting
\(\Box\varphi=\neg\Diamond\neg\varphi\) gives the normal rules
and local-premise convention of
Definition~\ref{def:p10-normal-modal-algebra-logic}; the proof below
verifies this comparison.
\index{normal modal logic!finitary signature}\index{polyadic modal logic}
\end{hbdefinition}

\begin{hbdefinition}{def:p10-polyadic-modal-algebra}{Algebraic interpretation and the finitary Lindenbaum algebra}
In a Boolean algebra \(\mathbf A\) with \(\Lambda\)-operators,
a valuation \(v:X\to A\) extends recursively to
\(\llbracket\varphi\rrbracket_v^{\mathbf A}\), interpreting each
\(f\) by \(f^A\).  For a nullary formula its value is \(f^A()\).
Write \(\mathbf A\models\varphi\) when every valuation gives value
\(1_A\), and put
\[
 \operatorname{Alg}_\Lambda(L)
   =\{\mathbf A:\mathbf A\text{ is a }\Lambda\text{-BAO and }
                     \mathbf A\models\lambda\text{ for every }\lambda\in L\}.
\]
Trivial algebras are included.
Define
\[
 \varphi\equiv_L\psi
       \quad\Longleftrightarrow\quad\vdash_L\varphi\leftrightarrow\psi.
\]
The next theorem proves this is a congruence.  Its quotient
\(F_{\Lambda,L}=\operatorname{Fm}_\Lambda(X)/{\equiv_L}\) has the
usual Boolean operations on classes and
\[
 f^{F_{\Lambda,L}}([\varphi_1]_L,\ldots,[\varphi_n]_L)
      =[f(\varphi_1,\ldots,\varphi_n)]_L,
 \qquad f^{F_{\Lambda,L}}()=[f()]_L\quad(n_f=0).
\]
\end{hbdefinition}

\begin{hbtheorem}{thm:p10-polyadic-algebra-completeness}{proved=annex:p10.polyadic.1}{Algebraic completeness for a finitary normal logic}
For every normal \(\Lambda\)-logic \(L\), the relation
\(\equiv_L\) is a congruence and
\(F_{\Lambda,L}\in\operatorname{Alg}_\Lambda(L)\).
For every formula \(\varphi\),
\[
 \vdash_L\varphi
 \quad\Longleftrightarrow\quad [\varphi]_L=1
 \quad\Longleftrightarrow\quad F_{\Lambda,L}\models\varphi
 \quad\Longleftrightarrow\quad
                  \operatorname{Alg}_\Lambda(L)\models\varphi.
\]
A non-theorem is refuted by \(p\mapsto[p]_L\) in this quotient.
If \(L\) is inconsistent, it contains every formula, its quotient
is trivial, and \(\operatorname{Alg}_\Lambda(L)\) consists exactly
of the trivial \(\Lambda\)-BAOs.
\end{hbtheorem}

This is the full finitary normal case of
\cite[preprint Definition~3.11, p.~9, and Theorem~3.13,
pp.~10--11]{Venema2005Preprint}.
The proof uses representatives only for the finitely many variables
in the formula under consideration, even when \(\Lambda\) is infinite.

\begin{hbdefinition}{def:p10-polyadic-canonical-frame}{The descriptive canonical frame}
Assume BPIT and let \(L\) be consistent.
Write \(\operatorname{Ult}(A)\) for the ultrafilters of the Boolean
reduct of \(A\).
On the ultrafilters of
the Boolean reduct of \(F_{\Lambda,L}\), put
\[
\begin{gathered}
 W_L=\operatorname{Ult}(F_{\Lambda,L}),\qquad
 \widehat a=\{U\in W_L:a\in U\},\qquad
 \mathcal A_L=\{\widehat a:a\in F_{\Lambda,L}\},\\
 R_f(U,U_1,\ldots,U_n)
 \quad\Longleftrightarrow\quad
 \forall a_1\in U_1\cdots\forall a_n\in U_n\quad
                f^{F_{\Lambda,L}}(a_1,\ldots,a_n)\in U.
\end{gathered}
\]
For \(n_f=0\), this means \(R_f(U)\) iff \([f()]_L\in U\).
Write \(\mathcal D_L=(W_L,(R_f),\mathcal A_L)\).
Its canonical valuation is \(v_0(p)=\widehat{[p]_L}\).
These are the constructions of
\cite[preprint Definition~5.12, p.~19, and Definition~5.24,
p.~22]{Venema2005Preprint}.

In any general \(\Lambda\)-frame, valuations take their variable
values in the admissible algebra.  Boolean satisfaction has its usual
clauses, and
\[
 \mathcal G,v,w\models f(\varphi_1,\ldots,\varphi_n)
 \quad\Longleftrightarrow\quad
 \exists w_1,\ldots,w_n\quad
 \left(R_f(w,w_1,\ldots,w_n)\ \land\
       \bigwedge_{i=1}^n\mathcal G,v,w_i\models\varphi_i\right).
\]
For \(n=0\), the right side is simply \(R_f(w)\).
Validity quantifies over every admissible valuation and every point.
\end{hbdefinition}

\begin{hbtheorem}{thm:p10-polyadic-canonical-truth}{proved=annex:p10.polyadic.2}{Descriptive representation and canonical truth}
Assume BPIT.  For consistent normal \(L\), the general frame
\(\mathcal D_L\) is nonempty and descriptive.
It validates every member of \(L\) under every admissible valuation,
and
\[
 \mathcal D_L,v_0,U\models\varphi
       \quad\Longleftrightarrow\quad [\varphi]_L\in U.
\]
Consequently \(\mathcal D_L\models\varphi\) iff \(\varphi\in L\).
\end{hbtheorem}

\begin{hbtheorem}{thm:p10-polyadic-local-completeness}{proved=annex:p10.polyadic.3}{Strong local completeness over descriptive general frames}
Assume BPIT, and let \(\operatorname{DGen}_\Lambda(L)\) consist of
all nonempty descriptive general \(\Lambda\)-frames validating
every member of \(L\) under every admissible valuation.
For every set \(\Gamma\subseteq\operatorname{Fm}_\Lambda(X)\) and
formula \(\varphi\),
\[
 \Gamma\vdash_L\varphi
       \quad\Longleftrightarrow\quad
 \Gamma\models^\ell_{\operatorname{DGen}_\Lambda(L)}\varphi.
\]
The right side means that, at every point of every such frame and
under every admissible valuation, truth of all premises at that
point implies truth of the conclusion there.
For consistent \(L\), every failure of derivability is witnessed
in \(\mathcal D_L\) under \(v_0\).
For inconsistent \(L\), the frame class is empty and both sides
hold for every \(\Gamma,\varphi\).
The same equivalence holds if all nonempty general frames validating
\(L\) are used instead of only the descriptive ones.
\end{hbtheorem}

The algebraic statement of Venema's Theorem~3.13 is a theoremhood
result.  The strong local result here uses the finite deduction
property and a Boolean filter containing the premise values and
the complement of the conclusion.  The general-frame representation
is descriptive by Theorem~\ref{thm:p10-bao-descriptive-representation},
corresponding to \cite[preprint Theorem~5.28, p.~23]{Venema2005Preprint}.
No claim that the underlying ordinary frame validates \(L\) under
all subset valuations is needed.
